\documentclass[10pt,a4paper]{amsart}
\usepackage[margin=19mm]{geometry}
\usepackage{amsmath,amssymb,mathtools}
\usepackage[scr=rsfs,cal=euler]{mathalfa}
\usepackage{xfrac}
\usepackage[unicode,hidelinks,bookmarks=false]{hyperref}
\newcommand{\ie}{i.e.\ }
\newcommand{\cf}{cf.\ }
\newcommand{\resp}{respectively\ }
\newcommand{\Subsection}{\subsection}
\providecommand{\nobreakdash}{\nobreak\hbox{-}}
\setlength{\emergencystretch}{2em}
\multlinegap0pt
\usepackage{bm}
\usepackage{bbm}
\usepackage{dsfont}
\usepackage{xspace}
\usepackage{cancel}
\usepackage{mathtools}
\usepackage{amsthm}
\makeatletter
\@ifundefined{theorem}{\newtheorem{theorem}{Theorem}}{}
\@ifundefined{lemma}{\newtheorem{lemma}[theorem]{Lemma}}{}
\@ifundefined{proposition}{\newtheorem{proposition}[theorem]{Proposition}}{}
\makeatother
\newcommand\numberthis{\addtocounter{equation}{1}\tag{\theequation}}
\newcommand{\cM}{\mathcal{M}}
\title[The second moment of cubic Dirichlet L-functions over functi]{The second moment of cubic Dirichlet L-functions over function fields}
\author{Shivani Goel, Anwesh Ray}
\date{Source: arXiv:2505.12015v2 (CC BY 4.0)}
\begin{document}
\maketitle

\noindent\emph{Source and licence.} The second moment of cubic Dirichlet L-functions over function fields. Version arXiv:2505.12015v2 (CC BY 4.0), distributed under the Creative Commons Attribution 4.0 International licence, as recorded in the arXiv licence metadata for this exact version and shown on the abstract page https://arxiv.org/abs/2505.12015v2. Full rights notice: licenses/arxiv-2505.12015-CC-BY-4.0.txt.\par\smallskip

\noindent\textit{Editorial scope.} Counted unit: the Proposition whose source label is \texttt{approx funct eqn} in the article (source0.tex lines 465--471, source label \texttt{approx funct eqn}), delivered together with the article's own complete original proof of it (source0.tex lines 473--527, one \texttt{proof} environment of 55 lines). No mathematics is written, repaired or re-derived: statement and proof are the article's own text line for line. Required context retained uncounted: functional equation lemma (lines 445--457). Every definition, display and label of those lines is retained unchanged. Omitted from this package and not redistributed: 1--444, 458--464, 472--472, 528--1139. Nothing inside a retained excerpt is omitted or altered: every retained body is delivered exactly as the article's lines are written, so no omission receipt of any kind accompanies this package. Citations: the retained excerpts carry no citation command at all, so no bibliography entry of the article is transcribed and no reference is redistributed. Reference adaptation: none was needed and none was made; every reference inside the delivered unit resolves inside it, so no unresolved reference or citation is produced. Printed slips kept literal: no printed word, symbol, hypothesis or number of the counted statement or of its proof was altered. Layout and class: the article's own class is \texttt{amsart} with packages including geometry, amscd, amsmath, amsxtra, amsthm, amssymb, stmaryrd, xr, mathrsfs, mathtools, enumerate, commath, comment, enumitem; the wrapper is the lane's pinned \texttt{amsart} editorial wrapper at 10pt on A4, which restates the theorem-like environments the retained text needs and adds the article's own macro definitions that the retained text needs (\texttt{\textbackslash cM}, \texttt{\textbackslash numberthis}), with extra wrapper packages (bm, bbm, dsfont, xspace, cancel, mathtools). The delivered numbering is the unit's own. Assets: no retained span contains a figure environment, a graphics-inclusion command, a PDF-page inclusion, a table or a TikZ picture; no image file is copied into this package, the package declares no asset dependency and no image file is redistributed with it. Licence: version arXiv:2505.12015v2 of this article is distributed under the Creative Commons Attribution 4.0 International licence, as recorded in the arXiv licence metadata for this exact version and shown on the abstract page https://arxiv.org/abs/2505.12015v2. A full rights notice is delivered at licenses/arxiv-2505.12015-CC-BY-4.0.txt. Characters: every retained excerpt is pure ASCII, so the delivered engine, XeLaTeX with its default Latin Modern text font, reports no missing character.
\begin{lemma} \label{functional equation lemma}
With respect to notation above, the following assertions hold.
\begin{enumerate}
    \item If $\chi$ is odd, we have that
    \begin{equation}\label{FE odd case}
\mathcal{L}_{q}(u, \chi) = \omega(\chi) (\sqrt{q} u)^{g} \mathcal{L}_{q} \left(\frac{1}{q u}, \bar{\chi} \right).
\end{equation}
\item On the other hand, suppose that $\chi$ is even, then 
\begin{equation}\label{FE even case}
\mathcal{L}_{q}(u, \chi) = \omega(\chi) (\sqrt{q} u)^{g} \frac{1 - u}{1 - \frac{1}{q u}} \mathcal{L}_{q} \left(\frac{1}{q u}, \bar{\chi} \right).
\end{equation}
\end{enumerate}
\end{lemma}

\begin{proposition}[Approximate functional equation]\label{approx funct eqn}
    Let $\chi$ be a nontrivial primitive cubic character of modulus $h$. If $\chi$ is odd, then
    \[L(1/2, \chi)^k=\sum_{f\in \cM_{\leq A}}\frac{\chi(f)d_k(f)}{\sqrt{|f|}}+ \omega(\chi)^k\sum_{f\in \cM_{\leq kg-A-1}}  \frac{\bar{\chi}(f) d_{k}(f)}{\sqrt{|f|}}.\]
    On the other hand, if $\chi$ is even, 
    \[\begin{split}L(1/2, \chi)^k=&(1-q^{-1/2})^k \sum_{i=0}^{A}\binom{k+i-1}{i}\sum_{f\in \cM_{\leq A-i}}\frac{\chi(f)d_k(f)}{\sqrt{|f|}} \\ + & (1-q^{-1/2})^k \omega(\chi)^k \sum_{i=0}^{kg-A-1}\binom{k+i-1}{i}\sum_{f\in \cM_{\leq kg-A-1-i}}  \frac{\bar{\chi}(f) d_{k}(f)}{\sqrt{|f|}}.
    \end{split}\]
\end{proposition}

\begin{proof}
First we assume that $\chi$ is odd, then by part (1) of Lemma \ref{functional equation lemma}, we have that
$$
\mathcal{L}_{q}(u, \chi)^k=\omega(\chi)^k(\sqrt{q} u)^{kg} \mathcal{L}_{q}\left(\frac{1}{q u}, \bar{\chi}\right)^k .
$$  Writing
$
\mathcal{L}_q\left(u, \chi\right)^{k}=\sum_{n=0}^{k g} b_n(\chi) u^{n}
$,
where $b_n(\chi)=\sum_{f \in \mathcal{M}_{n}} \chi(f) d_{k}(f)$, it follows that
$$
\sum_{n=0}^{k g} b_n(\chi) u^{n}=\omega(\chi)^k (\sqrt{q} u)^{kg} \sum_{n=0}^{k g} b_n(\bar{\chi}) \left(\frac{1}{qu}\right)^n
=\sum_{n=0}^{kg} b_n(\bar{\chi}) \omega(\chi)^k q^{kg/2-n} u^{kg-n}.$$
Therefore, we have that \[b_n(\chi)=b_{kg-n}(\bar{\chi})\omega(\chi)^k q^{n-kg/2},\] i.e., 
\[\sum_{f \in \mathcal{M}_{n}} \chi(f) d_{k}(f)=\sum_{f \in \mathcal{M}_{kg-n}} \omega(\chi)^k q^{n-kg/2}\bar{\chi}(f) d_{k}(f).\]
We find that:
\[\begin{split}\mathcal{L}_{q}(u, \chi)^k=& \sum_{n=0}^{A} \left(\sum_{f \in \mathcal{M}_{n}} \chi(f)d_k(f)\right) u^n+\sum_{n=A+1}^{kg} \left(\sum_{f \in \mathcal{M}_{n}} \chi(f)d_k(f)\right) u^n\\
= & \sum_{n=0}^{A} \left(\sum_{f \in \mathcal{M}_{n}} \chi(f)d_k(f)\right) u^n+ \sum_{n=A+1}^{kg} \left(\sum_{f \in \mathcal{M}_{kg-n}} \omega(\chi)^k q^{n-kg/2}\bar{\chi}(f) d_{k}(f)\right)u^n \\
= & \sum_{n=0}^{A} \left(\sum_{f \in \mathcal{M}_{n}} \chi(f)d_k(f)\right) u^n+ q^{kg/2}u^{kg}\omega(\chi)^k\sum_{n=0}^{kg-A-1} \left(\sum_{f \in \mathcal{M}_{n}} \bar{\chi}(f)  d_{k}(f)\right)(qu)^{-n}.
\end{split}\]
Plug in $u=q^{-1/2}$ to the above expression to get that 
\[L(1/2, \chi)^k=\sum_{f\in \cM_{\leq A}}\frac{\chi(f)d_k(f)}{\sqrt{|f|}}+ \omega(\chi)^k\sum_{f\in \cM_{\leq kg-A-1}}  \frac{\bar{\chi}(f) d_{k}(f)}{\sqrt{|f|}}.\] This proves part (1) of the Proposition.
\par Next we prove part (2) and assume that $\chi$ is even. In this case, we write
$$
F_\chi(u):=\left(\frac{\mathcal{L}_q\left(u, \chi\right)}{(1-u)}\right)^{k}=\sum_{n=0}^{k g} b_n(\chi) u^{n}.
$$
By the functional equation in part (2) of Lemma \ref{functional equation lemma}, 
\[F_\chi(u)=\omega(\chi)^k (\sqrt{q} u)^{kg}F_{\bar{\chi}}\left(\frac{1}{qu}\right)\]
and thus, $b_n(\chi)=b_{kg-n}(\bar{\chi})\omega(\chi)^k q^{n-kg/2}$.
It follows that 
\begin{align*}
    F_\chi(u)&= \sum_{n=0}^{A} b_n(\chi) u^{n}+\sum_{n=A+1}^{kg} b_n(\chi) u^{n}\\
    &=\sum_{n=0}^{A} b_n(\chi) u^{n}+\sum_{n=A+1}^{kg}b_{kg-n}(\bar{\chi})\omega(\chi)^k q^{n-kg/2}u^n\\
    &=\sum_{n=0}^{A} b_n(\chi) u^{n}+\omega(\chi)^k q^{kg/2}u^{kg}\sum_{n=0}^{kg-A-1}b_n(\bar \chi)(qu)^{-n}.\numberthis\label{functional equation}
\end{align*}
Now, if we  write
$$
\mathcal{L}_q\left(u, \chi\right)^{k}=\sum_{n=0}^{k g+k} a_n(\chi) u^{n}
$$
where $a_n(\chi)=\sum_{f \in \mathcal{M}_{n}} \chi(f) d_{k}(f)$.  %This gives
%\[a_n(\chi)=\sum_{j=0}^k \binom{k}{j}(-1)^j b_{n-j}(\chi)\]
%for $n=0,\ldots, kg$ provided that $b_n=0$ for $n<0$. and \[a_{kg+m}(\chi)=\sum_{j=m}^k \binom{k}{j}(-1)^j b_{kg+m-j}((\chi)\] for $m=0,\ldots, k.$
Note that $(1-u)^{-k}\mathcal{L}_q\left(u, \chi\right)^{k}=F_\chi(u)$. Using the Taylor expansion for $(1-u)^{-k}$ gives us the following relation: 
\begin{equation}\label{binomial inversion}
    b_n(\chi)=\sum_{i=0}^{n} \binom{k+i-1}{i} a_{n-i}(\chi).
\end{equation}
 Noting that $\mathcal{L}_q\left(u, \chi\right)^{k}=(1-u)^k\sum_{n=0}^{kg}b_n u^n$, we deduce from the functional equation \eqref{functional equation}:
\begin{align*}
    \mathcal{L}_q\left(u, \chi\right)^{k}&=(1-u)^k\left(\sum_{n=0}^{A} b_n(\chi) u^{n}+\omega(\chi)^k q^{kg/2}u^{kg}\sum_{n=0}^{kg-A-1}b_n(\bar \chi)(qu)^{-n}\right).
\end{align*}
Setting $u=q^{-1/2}$, we arrive at the following:
{\small\begin{align*}&\mathcal{L}_q\left(1/\sqrt{q}, \chi\right)^{k}\\
    & =(1-q^{-1/2})^k\left(\sum_{n=0}^{A} b_n(\chi) q^{-n/2}+\omega(\chi)^k \sum_{n=0}^{kg-A-1}b_n(\bar \chi)q^{-n/2}\right)\\&=(1-q^{-1/2})^k\sum_{i=0}^{A}\binom{k+i-1}{i}\left(\sum_{n=i}^{A} a_{n-i}(\chi) q^{-n/2}+\omega(\chi)^k \sum_{n=i}^{kg-A-1}a_{n-i}(\bar \chi)q^{-n/2}\right)\\&
 =  (1-q^{-1/2})^k \left(\sum_{i=0}^{A}\binom{k+i-1}{i}\sum_{f\in \cM_{\leq A-i}}\frac{\chi(f)d_k(f)}{\sqrt{|f|}}+ \sum_{i=0}^{kg-A-1}\binom{k+i-1}{i}\omega(\chi)^k\sum_{f\in \cM_{\leq kg-A-1-i}}  \frac{\bar{\chi}(f) d_{k}(f)}{\sqrt{|f|}}\right).
\end{align*}}
\end{proof}

\end{document}
